% Pursuing Stacks (cotes 134-2 to 134-8): archive cote 134-3, site page 2 (PDF page 3), typescript page 2.
% typescript. Draft compared: ps2.tex:55-115.
% AI-transcribed by Claude Opus 5.5 (claude-opus-5-5, adaptive thinking, effort high), Message Batches API, batch ps1 (msgbatch_01Rv5hCaNTtUt4JHx1yqoHPL), 2026-11-12; not yet checked by eye.
\scanpage{134-3/2}
\margin{2}
\mspage{2}

topological intuition,  by the people \inserted{primarily} concerned,  namely the topologists.
T\over{h}{h}e need of eliminating wildness phenomena has of course been felt repea-%
tedly,  and (by lack of anything better maybe,  or rather by lack of any
attempt of a systematic reflection on what \uline{was} needed) it was supposed
to be met by the notion of piecewise linear structures.  This however
was falling from one extreme into another  -  from a structure species
with vastly too many maps between "spaces",  like a coat vastly too wide
and floating around in a million wild wrinkles,  to one with so few
(not even a quadratic map from $\underline{R}$ to $\underline{R}$ is allowed !) that it feels like
too narrow a coat,  bursting \inserted{apart} on all edges and ends.  The main defect
here,  technically speaking,  seems to me the fact that \deleted{numerical}\inserted{$\mathbb{R}$-valued} piece-%
wise linear functions are not stable under multiplication,  and as a geo-%
metric consequence of this,  that when contracting a compact p.l. sub-%
space of a (compact, say) piecewise linear space into one point,  we do
not get on the quotient \xout{\unsure{structurexx}} \hand{space} a natural p.l. structure.  This
alone should have sufficed,  one might think,  to eliminate the piecewi-%
se linear structure species as a reasonable candidate for \xout{eliminating}
``doing topology'' wi\over{th}{th}out wildness empidements  -  but strangely enough,
it seems to be hanging around till this very day !

\margin{(15.)}15.  But my aim here is not to give an outline of foundations  on ``tame
topology'',  but rather to fill some foundational gap\deleted{s} in homotopy theory\hand{,}
more specifically in homotopical algebra.
\margin{pas à la ligne}The relative success \deleted{here} in th\over{e}{e} \deleted{\unsure{x topic,}}  homotopical approach to topo-%
logy is probably closely tied to the well known fact (\over{b}{B}rouwers starting
point as a matter of fact,  when he introduced the systematics of trian-%
gulations)  that every continuous map between triangulated spaces can be
approximated by simplicial maps.  This gave rise,  rather naturally,  to
the hope expressed in the \hand{``}Hauptvermutung\hand{''}  -  that two homeomorphic trian%
gulated spaces admit isomorphic subdivisions,  a hope that finally proved
a delusion.  With a distance of two or three generations,  I would
comment on this by saying that this negative result was the one to be
expected,  once it has become clear that neither of the two structure
species one was comparing,  namely topological spaces and triangulated
spaces,  was adequate for expressing what one is really after  -  namely
an accurate mathematical transcription,in terms of ``spaces'' of some kind
or other, of some vast and deep and misty and ever transforming \unsure{mase|mass} of
intuitions in our psyche,  which we are referring to as ``topological''
\margin{\unsure{$\Pi$}}intuition\deleted{s}\hand{$(^1)$}. There \uline{is}  something positive though,  definitely,  which
can be viewed as an extremely weakened version of the Hauptvermutung,
namely the fact that topological spaces on the one hand,  and semi-simp-%
plicial sets on the other,  give rise,  by a suitable ``localization''
